\documentclass[10pt]{article}

\usepackage[utf8]{inputenc}

\usepackage[T1]{fontenc}

\usepackage{amsmath}

\usepackage{amsfonts}

\usepackage{amssymb}

\usepackage[version=4]{mhchem}

\usepackage{stmaryrd}

\usepackage{graphicx}

\usepackage[export]{adjustbox}

\graphicspath{ {./images/} }

\usepackage{bbold}



\begin{document}

Если граница $\partial G$ - гладкое $(n-1)$-мерное многообразие, функция $\varphi$ гладкая и выполнено условие согласования на $\{t=0, x \in \partial G\}$, то П. к. з. однозначно разрешима.



Лит.: [1] Б и ц а д з е А. В., Некоторые классы уравнений в частных производных, М., 1981; [2] Бе р с Л., Д ж о н Ф., II е х т е p M., Уравнения с частными производными, пер.



\includegraphics[max width=\textwidth, center]{2023_07_08_e4ab63d2527305f3d3a2g-1(3)}

производными, пер. с англ., М., 1964; [4] М и р а н д а К., Уравнения с частными производными әллиптического типа, пер.

с итал., М., 1957; [5] П е т р о в с к и й И. Г. Лекции об уравнениях с частными производными, 3 изд., М., 1961; [6] Б иц а д з е А. В., Уравнекия математической физики, М., 1976; торы с частными производными, пер. с англ., М., 1965



A. ̈. С. Солдатов,



ПЕРВИЧНОЕ КОЛЬЦО - кольцо $R$, в к-ром произведение любых двусторонних идеалов $P$ и $Q$ равно нулевому идеалу в том и только в том случае, когда либо $P$, либо $Q$ является нулевым идеалом. Другими словами, идеалы П. к. по умножению образуют полугруппу без делителей нуля. Кольцо $R$ является П. к. тогда и только тогда, когда правый (левый) аннулятор любого его ненулевого правого (соответственно левого) идеала равен $(0)$, а также тогда и только тогда, когда $a R b \neq 0$ для любых ненулевых $a, b \in R$. Центр П. к. является областью челостности. Любое примитивное кольцо первично. Если кольцо $R$ не соцержит ненулевых нильидеалов, то $R$ - подпрямая сумма первичных колец. Класс П. к. играет важную роль в теории радикалов колец (см. [1]).



Существует следующее обобщение понятия П. к. Кольцо $R$ наз. п о л у п е р в и ч ны м, если оно не имеет ненулевых нилыпотентых идеалов.



\includegraphics[max width=\textwidth, center]{2023_07_08_e4ab63d2527305f3d3a2g-1}

Ю. М., Радикалы алгебр и структурная теория, М., $1979 ;$ [2] [3] X е р с т е й н Й., Некоммутативные кольца, пер. с англ., М., 1972.



ПЕРВИЧНЫИ ИДЕАЛ - такой дв усторонний идеал $I$ колыца $A$, что из включения $P Q \subseteq I$ для любых двусторонних идеалов $P$ и $Q$ кольца $\bar{A}$ следует, что либо $P \subseteq I$, либо $Q \subseteq I$. Для первичности идеала $I$ кольца $R$ необходимо и достаточно, чтобы множество $R \backslash I$ было m-с и с т е м о й, т. е. чтобы для любых $a, b \in R \backslash I$ существовал $x \in R$ такой, что $a x b \in R \backslash I$. Џдеал $I$ кольца $A$ первичен тогда и только тогда, когда факторколыцо по нему является первичныл кольцом.



H. А. жевлаков.



ПЕРВООБРАЗНАЯ, П е р о о б р а з в я (приМ и т и в н а я) ф у кция, для конечной функции $f(x)$ - такая функция $F(x)$, что всюду $F^{\prime}(x)=f(x)$. Это определение является наиболее распространенным, но встречаются и другие, в к-рых ослаблены требования сущөствования всюду конечной $F^{\prime}$ и вышолнения всюду равенства $F^{\prime}=f$; иногда в определении используют обобщения производной. Большинство теорем 0 П. касается их существования, нахождения и единственности. Достаточным условием для существования П. у заданной на отрезке функции $f$ является непрерывность $f$; необходимыми условиями являются принадлежность функции $f$ первому Бәра классу и выполнение для нее Дарбу свойства. У заданной на отревке функции любые две П. отличаются на постоянную. Задачу вахождения $F$ по $F^{\prime}$ для нешрерывных $F^{\prime}$ репает $P u-$ мана интеграл, для ограниченных $F^{\prime}$ - Лебега интеграл, для любой $F^{\prime}-$ узкий (а тем болеө широкий) Данжуа интеграл и Перрона интеграл.



Лum.: [1] Ғ у д р я в ц ө в Л. Д., Курс математического анализа, т. 1, М., 1981; [2] Н и к о л ь с к и й С. М., Курс матеанализа, т. 1, М., $1981 ;$ й н и к о л ь с к и й С. М., Курс мате-

матического анализа, 2 изд., т. 1, М., 1975. Т. П. Лyкашенко.



ПЕРВООБРАЗНЫИИ КОРЕНБ - 1) П. к., П р У м ит и в ы й к о р е н ь, из единицы в поле $K$ степени $m$ - әлемент $\zeta$ поля $K$ такой, что $\zeta^{m}=1$ и $\zeta^{r} \neq 1$ для любого натурального $r<m$. Элемент $\zeta$ порождает циклич. группу $\mu(m)$ корней из единицы порядка $m$. Если в поле $K$ существует П. к. степени $m$, то $m$ взаимно просто с характеристикой поля $K$. Алгебраически замкнутое поле содержпт П. к. любой степени взаимно простой с характеристикой поля. Если $\zeta-$ П. к. степени $n$, то для любого $k$ взаимно простого с $n$ элемент $\zeta^{k}$ также является П. к. Число всех П. к. степени $m$ равно значению функции Эйлера $\varphi(m)$.



В поле комплексных чисел П. к. степени $m$ имеют вид



\[

\cos \frac{2 \pi k}{m}+i \sin \frac{2 \pi k}{m}

\]



где $0<k<m$ и $\boldsymbol{k}$ взаимно просто с $m$.



\begin{enumerate}

  \setcounter{enumi}{1}

  \item П. к. І о м о д у л ю $m$ - целое число $g$ такое, что

\end{enumerate}



\begin{center}

\includegraphics[max width=\textwidth]{2023_07_08_e4ab63d2527305f3d3a2g-1(1)}

\end{center}



при $1 \leqslant \gamma \leqslant \varphi(m)-1$, где $\varphi(m)$ - функция Эйлера. Для П. к. $g$ его степени $g^{0}=1, g$, .., $g \varphi(m)-1$ несравнимы между собой по модулю $m$ и образуют приведенную систему вычетов по модулю $m$. Таким образом, для каждого числа $a$, взаимно простого с $m$, найдется показатель $\mathcal{}(0 \leqslant \gamma \leqslant \varphi(m)-1)$, для к-рого $a \equiv g \gamma(\bmod m)$.



П. к. существуют не для всех модулей, а только для модулей $m$ вида $m=2,4, p \alpha, 2 p \alpha$, где $p>2-$ простое число. В этих случаях мультипликативные группь приведенных классов вычетов по модулю $m$ устроены наиболее просто: они являются циклич. группами порядка $\varphi(m)$. С понятием П. к. по модулю $m$ тесно связано понятие индекса числа по модулю $\mathrm{m}$.



П. к. для простых модулей $p$ были введены Л. Эйлером (L. Euler), но существование П. к. для любых простых модулей $p$ было доказано лишь К. Гауссом (C. Gauss, 1801).



Лит.: Л е н г С., Алгебра, пер. с англ., М., 1968; [2] Гау с с К. Ф., Труды по теории чисел, пер. с лат. и нем., М., 1959 [3] В и но г р а д о в И. М., Основы теории чисел, 8 изд., М.,

1972 .

Л. В. Кузвмин, С. А. Степанов.



ПЕРВЫИ ИНТЕГРАЛ ди фф е ренци а льного ур а внения н-отличная от постоянной и непрерывно дифференцируемая Функция, производная к-рой вдоль решений данного уравнения тождественно равна нулю. Для скалярного уравнения



\[

y^{\prime}=f(x, y)

\]



П. и. есть функция $F(x, y)$, находящаяся в левой части общего решения $F(x, y)=C$, где $C$ - произвольная постоянная. Таким образом, $F(x, y)$ удовлетворяет линейному уравнению



\[

\frac{\partial F(x, y)}{\partial x}+\frac{\partial F(x, y)}{\partial y} f(x, y)=0

\]



с частными производными 1-го порядка. П. п. может не существовать во всей области задания уравнения (*), однако в малой окрестности точки, в к-рой функция $f(x, y)$ непрерывно дифференцируема, он всегда существует. П. и. определяется не единственным образом. Так, для уравнения $y^{\prime}=-x / y$ П. и. является как функция $x^{2}+y^{2}$, так, напр., и функция $\exp \left(x^{2}+y^{2}\right)$.



Знание П. и. нормальной системы



\[

\dot{x}=f(x, t), \quad x \in \mathbb{R}^{n},

\]



позволяет понизить порядок этой системы на единицу, а отыскание $n$ функционально независимых П. и. равносильно отысканию общего решения в неявном виде.



\includegraphics[max width=\textwidth, center]{2023_07_08_e4ab63d2527305f3d3a2g-1(2)}

висимые П. и., то всякий другой П. и. $F(x, t)$ можно представить в виде



\[

F(x, t)=\Phi\left(F_{1}(x, t), \ldots, F_{n}(x, t)\right),

\]



где Ф - нек-рая дифференцируемая функция.



Лит.: [1] П о н т р я г и н Л. С., Обыкновенные дифференциальные уравнения, 4 изд., М., 1974.



H. H. Лaวuc.





\end{document}